\section{Основные этапы развития нелинейной оптики}\label{sec:chFiber3/fiberHist}

Впервые идеи нелинейной оптики были высказаны советским учёным С.И.~Вавиловым в работе <<Микроструктура света>> \cite{vavilov1950}. Он утверждал, что нелинейность по величине световой мощности должна наблюдаться в отношении не только абсорбции, но и других оптических свойств среды: двулучепреломления, дихроизма, вращательной способности и т.д. Вавилов указал лишь на часть физических явлений, приводящих к нелинейности среды по световому полю (изменение числа атомов или молекул, способных поглощать свет), и, как следствие, ряд  нелинейных явлений, основанных на других физических механизмах, были предсказаны и продемонстрированы позднее, С. И. Вавилов по праву считается основоположником нелинейной оптики.	 

Впервые генерация второй гармоники (ГВГ) \cite{franken1961generation} наблюдалась спустя всего год после создания первого лазера в 1960 году \cite{maiman1960stimulated}. Несмотря на то, что для эффективного протекания процесса необходимо соблюдение условий т.н. фазового синхронизма, в данной работе эффект наблюдался в его отсутствие. Затем в 1962 году А. Армстронгом и Н. Бломбергеном  были предложены три способа достижения фазового синхронизма дискретного типа или квазисинхронизма \cite{armstrong1962interactions}, при котором периодически контролируемо изменялась разность фаз взаимодействующих волн, а американский учёный Джордмэйн предложил использовать для достижения этой цели двулучепреломляющие кристаллы \cite{giordmaine1962mixing}. Получение синхронизма с использованием двулучепреломления получило широкое применение, и данный тип синхронизма получил название традиционного. 

Последовавший за публикацией \cite{giordmaine1962mixing} бурный рост числа исследований по генерации второй гармоники на основе традиционного синхронизма привёл к получению уже в 1963 году преобразования с эффективностью 20\% ССІЛКА, а идеи Н.~Бломбергена с сотрудниками в силу технологических трудностей создания структур для реализации квазисинхронизма были отодвинуты на второй план. Лишь в 90-х годах XX века произошёл возврат к идеям квазисинхронизма на основе кристаллов с регулярной доменной структурой.


\fixme{Н-о процессы третьего порядка изучались как в кристаллах и стёклах, так и в оптических волокнах, причем в последних они эти процессы проявляются ярче в силу модовой структуры излучения и, как следствие, высоких интенсивностей излучения на больших длинах оптического волокна.}
\textbf{Комбинационное} или {Рамановское} рассеяние было открыто независимо советским учёным Мандельштамом в кристаллах и индийским учёным Раманом~--- в жидкостях, в 1928 году \textbf{ССЫЛКА}. Присуждение за это открытие Нобелевской премии \fixme{в каком году} только Раману до сих пор является \fixme{предметом взаимных упрёков.} В 1962 году \textbf{ССЫЛКА} с использованием \fixme{рубинового} лазера был открыт эффект вынужденного комбинационного рассеяния (ВКР).
\nomenclature{ВКР}{Вынужденное комбинационное рассеяние}

Рассеяние света на акустических фононах в твёрдых телах и жидкостях было предсказано в 1914 году \textbf{ССЫЛКА} и экспериментально продемонстрировано в 1930 году \textbf{ССЫЛКА}. Вынужденное же рассеяние лазерного излучения на таких фононах было обнаружено в 1964 году \textbf{ССЫЛКА} и носит название \textbf{вынужденного рассеяние Мандельштама-Бриллюэна (ВРМБ).}
\nomenclature{ВРМБ}{Вынужденное рассеяние Мандельштама-Бриллюэна}

Оба вышеназванных эффекта активно исследовались в волоконных световодах в 1972 г. ССЫЛКИ, поскольку к этому моменту уровень оптических потерь в световодах приблизился к фундаментальному ограничению, обусловленному процессом Рэлеевского рассеяния, и нелинейно-оптические эффекты третьего порядка стали проявляться особенно ярко.

Нелинейно-оптические эффекты третьего порядка, обусловленные зависимостью показателя преломления среды от величины светового поля и не зависящение от фазовых соотношений взаимодействующих волн~--- фазовая кросс-модуляция (ФКМ) и фазовая самомодуляция (ФСМ)~--- были предсказаны и экспериментально обнаружены в конце 1960-х--начале 1970-х годов и в 1970-х годах \textbf{ССЫЛКА} соответственно. Оказалось, что оптические импульсы высокой интенсивности изменяют показатель преломления среды, в которой они распространяются, что, в свою очередь, приводит к изменению спектрального состава самих импульсов.
\nomenclature{ФСМ}{Фазовая кросс-модуляция}
\nomenclature{ФКМ}{Фазовая самомодуляция}
\nomenclature{ОВ}{Оптическое волокно}
\nomenclature{н-о}{нелинейно-оптический}

Помимо упомянутых выше, существуют также н-о процессы третьего порядка, требующие фазового согласования взаимодействующих волн - выполнения условия фазового синхронизма. К ним относят ЧВС~--- четырёхволновое смешение. Этот процесс также был впервые теоретически предсказан, описан и обнаружен в конце 1960-х -- начале 1970-х годов \textbf{ССЫЛКА}. И здесь использование ОВ привело к активному изучению этого эффекта. В настоящее время ЧВС является одним из негативных факторов в системах оптической связи, но в то же время он может быть использован для создания новых источников света, генерации суперконтинуума, а также в квантовой оптике для генерации спутанных фотонов.